\section{METRによるAnthropicのAI開発ペース加速評価}
METRがAnthropic内部データに基づき、モデル能力向上ペースが過去に比べ約1.9倍（約2倍）加速していると評価した内容が観測窓内で議論された。AI支援による研究加速が主因とされ、従業員の自己報告生産性向上（最大12倍程度）も言及されている。 X上で複数の投稿がMETRの評価を要約・引用。 元の発表は観測窓外であり、本号は窓内の反応を記録する。
\dailyxsource{窓内の関連投稿}{https://x.com/TechElyra/status/2119352223662805308}

\section{NVIDIA Cosmos 3 Nanoオープンウェイトチェックポイント公開}
NVIDIAがCosmos 3 Nano（16Bパラメータ）ベースの4つのオープンウェイトチェックポイント（Transfer-Auto、Sim-Auto、Sim-Transfer、Sim-Bimanual）をHugging Faceで公開。運転シミュレーションや二腕ロボット向けのインタラクティブ世界モデルとして位置づけられ、OpenMDW 1.1ライセンスで商用利用可能とされている。 窓内の投稿が詳細を説明し、HFへのリンクや仕様を記載。 元の発表は観測窓外であり、本号は窓内の反応を記録する。
\dailyxsource{窓内の関連投稿}{https://x.com/Hesoyam676767/status/2119544238233055673}

\section{フロンティアモデル公開時期に関するコミュニティ予測}
Gemini 4 Carbonの公開接近、GPT-6.2 Astra/Sol、Claude Fable 5.5、Grok 4.8、DeepSeek V4.1 Pro、Qwen系などの「今週中」予測が複数投稿で共有された。公式発表は確認できず、過去のロールアウト間隔からの推測が中心。 複数の独立投稿が類似リストを挙げ、Googleの過去間隔を根拠にCarbon公開を予想。 これらはGrokが未確認と判定した投稿上の主張であり、確定した発表としては扱わない。
\dailyxsource{窓内の関連投稿}{https://x.com/daveeyyy_/status/2119527768161652914}

\section{Google DeepMind Genie 4世界モデルに関する議論}
Genie 4（リアルタイムインタラクティブ世界モデル、ネイティブオーディオ・テキストステアリング対応）に関する紹介・解説投稿が窓内に存在。発表自体は10月末。 窓内投稿が機能（リアルタイム変更、オーディオ生成）を説明。 元の発表は観測窓外であり、本号は窓内の反応を記録する。
\dailyxsource{窓内の関連投稿}{https://x.com/sajbersfera/status/2119548453374329222}

\section{Sunoに対するElvis権利者・Eminemパブリッシャー訴訟}
Elvis Presley権利者とEminemのパブリッシャー（Eight Mile Style）が、AI音楽生成Sunoに対し音声クローニング・著作権侵害を主張する訴訟を提起したと報じられ、X上で言及。 窓内投稿が訴訟を報じている。
\dailyxsource{窓内の関連投稿}{https://x.com/ababnews_com/status/2119550335362449677}

\section{ローカルコーディングエージェントツールTenazの更新}
ローカルモデル向けコーディングエージェントTenaz 0.6.0がDeepSeek HarnessサポートとACPエージェント対応を追加したと投稿。長時間実行の改善も報告。 投稿者が機能追加（DeepSeek Harness、キャッシュ改善など）を詳述。
\dailyxsource{窓内の関連投稿}{https://x.com/aartiles24/status/2119550069905014991}

\section{Hugging Faceトレンドオープンモデル（DeepSeek派生など）}
HFトレンドでDeepSeek-V4-Pro-Qwen3.5-9B-UNCENSORED派生などが上位にあり、ローカル実行可能なGGUFが共有された。NVIDIA Cosmosも言及。 投稿者がトレンド順位とファイルサイズを記載。 これらはGrokが未確認と判定した投稿上の主張であり、確定した発表としては扱わない。
\dailyxsource{窓内の関連投稿}{https://x.com/KrEarthMilkyway/status/2119549730267046238}

\section{OpenAI GPT-6 Cyberに関する言及}
OpenAIが脆弱性発見向けGPT-6 Cyberを発表・限定提供したとする投稿。Daybreak Red顧客向けなど具体的な能力主張あり。 投稿者主張のみ。公式一次資料は窓内で確認できず。 これらはGrokが未確認と判定した投稿上の主張であり、確定した発表としては扱わない。
\dailyxsource{窓内の関連投稿}{https://x.com/MoguriConsultan/status/2119550028230344714}

\section{AIデータセンター・電力インフラ動向}
GoogleがCrusoeの未完成データセンタープロジェクトを引き継いだとの報道や、Anthropicの電力関連契約などが窓内投稿で共有された。 投稿が複数ニュースをまとめている。 元の発表は観測窓外であり、本号は窓内の反応を記録する。
\dailyxsource{窓内の関連投稿}{https://x.com/KrEarthMilkyway/status/2119549730267046238}

\section{Anthropicの科学研究生産性分析言及}
Anthropicが科学者セッション分析でAIが研究時間の一定割合を担えるとするシナリオモデルを公開したと投稿で言及。 投稿者がリンク付きで要約。 これらはGrokが未確認と判定した投稿上の主張であり、確定した発表としては扱わない。
\dailyxsource{窓内の関連投稿}{https://x.com/KrEarthMilkyway/status/2119549730267046238}
